\uuid{4sBg}
\titre{Théorème de Bézout}
\auteur{bodin ; croizat ; sacre}
\organisation{exo7}
\licence{CC BY-NC-SA 4.0}
\datecreate{2022-03-08}
\niveau{L1}
\module{Algèbre}
\chapitre{Arithmétique dans Z}
\sousChapitre{Nombres premiers, nombres premiers entre eux}
\difficulte{2}
\theme{Théorème de Bézout}
\exo7qcmid{hilisit.arithmetique.020}

\contenu{
\begin{question}
\qclassification{103.04}
    Pour $a=453$ et $b=201$, l'algorithme d'Euclide (étendu) fournit des coefficients de Bézout $u$ et $v$ tels que $au+bv=\pgcd(a,b)$ avec :
    \begin{answers}
        \bad{$u=4$, $v=-9$, $\pgcd(a,b)=1$.}
        \bad{$u=-12$, $v=27$, $\pgcd(a,b) = 51$.}
        \bad{$u=1$, $v=-2$, $\pgcd(a,b)=51$}
        \good{$u=4$, $v=-9$, $\pgcd(a,b)=3$.}
    \end{answers}
    \begin{explanations}
        L'algorithme d'Euclide nous fournit :
        $$ 453 = 2 \times 201 + 51$$
        $$201 = 3 \times 51 + 48 $$
        $$ 51 = 1 \times 48 + 3$$
        $$48 = 16 \times 3 + 0 $$
        Ainsi on a $\pgcd(a,b) = 3$. En remontant cet algorithme, on obtient :
        $$ 3 = 51 - 48 = 201 \times (-1) + 51 \times 4 = 453 \times 4 + 201 \times (-9)$$
    \end{explanations}
\end{question}
}
